\begin{mvtable}
\mvtabcap{续表}
\begin{tabular}{ccccccc}
\toprule
编号 & $x_1$ & $x_2$ & $x_3$ & $x_4$ & $x_5$ & $x_6$\\
\midrule
30 & $-0.36$ & $-0.16$ & $0.15$ & $-0.38$ & $1.18$ & $-0.10$\\
31 & $-0.40$ & $-0.12$ & $-0.48$ & $-0.34$ & $0.30$ & $-0.20$\\
32 & $-0.60$ & $-0.40$ & $-0.20$ & $0.32$ & $0.50$ & $0.10$\\
33 & $-0.47$ & $-0.16$ & $-0.34$ & $-0.31$ & $0.85$ & $0.60$\\
34 & $-0.46$ & $-0.18$ & $0.16$ & $0.01$ & $0.60$ & $0.35$\\
35 & $-0.44$ & $-0.12$ & $-0.20$ & $-0.48$ & $1.40$ & $0.10$\\
36 & $-0.90$ & $-0.40$ & $0.75$ & $-0.31$ & $0.60$ & $-0.10$\\
37 & $-0.50$ & $-0.35$ & $0.84$ & $-0.52$ & $0.35$ & $-0.75$\\
38 & $-0.38$ & $0.08$ & $0.55$ & $-0.15$ & $0.80$ & $-0.10$\\
39 & $-0.60$ & $-0.35$ & $-0.35$ & $-0.34$ & $0.60$ & $0.85$\\
40 & $0.11$ & $0.24$ & $0.15$ & $0.40$ & $0.00$ & $-0.10$\\
41 & $0.05$ & $0.12$ & $0.85$ & $0.55$ & $1.65$ & $-0.10$\\
42 & $-0.85$ & $-0.65$ & $0.50$ & $0.35$ & $0.80$ & $-0.21$\\
43 & $-0.37$ & $-0.10$ & $-0.10$ & $-0.58$ & $1.85$ & $-0.11$\\
44 & $-0.11$ & $0.24$ & $0.75$ & $-0.10$ & $0.65$ & $-0.10$\\
45 & $-0.60$ & $-0.24$ & $0.13$ & $0.84$ & $0.85$ & $0.15$\\
46 & $-0.84$ & $-0.59$ & $0.05$ & $0.61$ & $1.00$ & $0.20$\\
47 & $-0.46$ & $-0.16$ & $0.37$ & $-0.15$ & $0.68$ & $0.25$\\
48 & $-0.56$ & $-0.35$ & $-0.10$ & $0.75$ & $0.45$ & $0.20$\\
49 & $-0.56$ & $-0.16$ & $0.37$ & $-0.25$ & $1.05$ & $0.15$\\
50 & $-0.25$ & $-0.12$ & $-0.05$ & $-0.20$ & $1.21$ & $0.10$\\
\bottomrule
\end{tabular}
\end{mvtable}

资料来源：数据经 Darek Ceglarek 许可使用.

\noindent 5.29\quad 参考练习 5.28 中的车体数据.这些数据都是作为目标值的离差测量的,所以用来检验均值向量为零的零假设是合适的.使用前 30 个数据,在水平 $\alpha=0.05$ 下,检验 $H_0$: $\boldsymbol{\mu}=\mathbf{0}$.

\noindent 5.30\quad 参考练习 3.18 中的能量消耗的数据.

(a) 对 4 种类型中的每一种能源的平均消耗,它们的平均总消耗量和汽油减去天然气的平均差,求大样本的 95\% 庞弗罗尼置信区间.

(b) 对 4 种类型中的每一种能源的平均消耗,它们的平均总消耗量和汽油减去天然气的平均差,求大样本的 95\% 联合 $T^2$ 区间.并与 (a) 中的结果进行比较.

\noindent 5.31\quad 参考练习 3.20 中的暴风雪数据.

(a) 做一些合适的变换之后,找出均值向量的 95\% 的置信域.

(b) 在同样的尺度下,求两个均值分量的 95\% 庞弗罗尼置信区间.

\begin{center}
\phantomsection\label{bib:5}
\textbf{参考文献}
\end{center}

\noindent 1. Agresti, A. \textit{Categorical Data Analysis} (2nd ed.), New York: John Wiley, 2002.

\noindent 2. Bacon-Sone, J., and W. K. Fung. ``A New Graphical Method for Detecting Single and Multiple Outliers in Univariate and Multivariate Data.'' \textit{Applied Statistics}, \textbf{36}, no. 2 (1987), 153--162.
